\documentclass[12pt]{article}
\usepackage[T1]{fontenc}
\usepackage[utf8]{inputenc}
\usepackage{lmodern}
\usepackage{textcomp}
\usepackage{enumitem}
\usepackage{geometry}
\geometry{margin=1in}
\begin{document}
\begin{center}
Answer Key - Machine Learning - Graduate Level Assessment 5 \
\small (Answers are ordered exactly as in the question paper.)
\end{center}

\section*{Multiple Choice}
\begin{enumerate}[label=\arabic*.]
\item \textbf{Answer:} B
\\ \textit{Options:} A) Increase the amount of training data.; B) Improve the optimisation algorithm being used for error minimisation.; C) Decrease the model complexity.; D) Reduce the noise in the training data.
\\ \textit{Note:} Correct option: B - Improve the optimisation algorithm being used for error minimisation.
\item \textbf{Answer:} D
\\ \textit{Options:} A) It can not be applied to non-differentiable functions.; B) It can not be applied to non-continuous functions.; C) It is hard to implement.; D) It runs reasonably slow for multiple linear regression.
\\ \textit{Note:} Correct option: D - It runs reasonably slow for multiple linear regression.
\item \textbf{Answer:} D
\\ \textit{Options:} A) It relates inputs to outputs.; B) It is used for prediction.; C) It may be used for interpretation.; D) It discovers causal relationships
\\ \textit{Note:} Correct option: D - It discovers causal relationships
\item \textbf{Answer:} A
\\ \textit{Options:} A) True, True; B) False, False; C) True, False; D) False, True
\\ \textit{Note:} Correct option: A - True, True
\item \textbf{Answer:} B
\\ \textit{Options:} A) The use of sampling with replacement as the sampling technique; B) The use of weak classifiers; C) The use of classification algorithms which are not prone to overfitting; D) The practice of validation performed on every classifier trained
\\ \textit{Note:} Correct option: B - The use of weak classifiers
\end{enumerate}

\section*{True / False}
\begin{enumerate}[label=\arabic*.]
\setcounter{enumi}{5}
\item \textbf{Answer:} True
\\ \textit{Options:} A) True; B) False
\\ \textit{Note:} LLM-generated true statement based on MCQ.
\item \textbf{Answer:} True
\\ \textit{Options:} A) True; B) False
\\ \textit{Note:} LLM-generated true statement based on MCQ.
\item \textbf{Answer:} True
\\ \textit{Options:} A) True; B) False
\\ \textit{Note:} LLM-generated true statement based on MCQ.
\item \textbf{Answer:} True
\\ \textit{Options:} A) True; B) False
\\ \textit{Note:} LLM-generated true statement based on MCQ.
\item \textbf{Answer:} True
\\ \textit{Options:} A) True; B) False
\\ \textit{Note:} LLM-generated true statement based on MCQ.
\end{enumerate}

\section*{Long Form Response}
\begin{enumerate}[label=\arabic*.]
\setcounter{enumi}{10}
\item \textbf{Answer:} def \_\_init\_\_(self, data): | def max\_height(node): | return 0 ; | return left\_height+1 | return right\_height+1
\\ \textit{Note:} Students should provide a detailed explanation referencing algorithm steps.
\item \textbf{Answer:} def count\_elim(num): | return count\_elim
\\ \textit{Note:} Students should provide a detailed explanation referencing algorithm steps.
\end{enumerate}

\vfill
\noindent\textit{End of Answer Key}
\end{document}